\section{Praktikum}
\label{sec:praktikum-bab4}

\begin{praktikumbox}
\vspace{1.1em}
\textbf{Judul:} Diagram blok 8086 dan penelusuran siklus instruksi.\\[0.3em]
\textbf{Sub-CPMK:} 1.2 \quad\textbullet\quad
\textbf{Alokasi waktu:} 100 menit \quad\textbullet\quad
\textbf{Bentuk:} individu
\end{praktikumbox}

\subsection{Tujuan Praktikum}

Setelah menyelesaikan praktikum ini, mahasiswa mampu menggambar
diagram blok BIU/EU, menjelaskan pembagian tugas kedua unit, serta
menelusuri perubahan \reg{IP}, \reg{AX}, dan isi \eng{queue} pada
setiap langkah eksekusi.

\subsection{Alat dan Bahan}

\begin{itemize}
    \item Komputer dengan emu8086 terpasang.
    \item DOSBox dengan TASM, TLINK, dan DEBUG.
    \item Berkas \file{siklus-dasar.asm} dari \file{code/asm/}.
    \item Kertas kosong, alat tulis, dan lembar jurnal praktikum.
\end{itemize}

\subsection{Langkah Kerja}

\begin{enumerate}
    \item \textbf{Menggambar diagram blok dari ingatan.}
    \begin{enumerate}
        \item Tutup halaman materi, lalu gambarkan kotak EU dan BIU
              beserta bus internal dan bus eksternal.
        \item Tuliskan komponen yang berada di dalam EU dan BIU,
              termasuk \eng{instruction queue} enam byte.
        \item Buka kembali Gambar~\ref{fig:blok-8086} dan tandai
              bagian yang salah letak atau terlewat.
    \end{enumerate}

    \item \textbf{Menelusuri program dengan emu8086.}
    \begin{enumerate}
        \item Buka \file{siklus-dasar.asm}, lalu pilih
              \file{Compile and Emulate}.
        \item Catat nilai awal \reg{CS}, \reg{IP}, \reg{AX}, dan
              \reg{DS}.
        \item Tekan \file{Single Step} satu kali. Catat perubahan
              \reg{IP}, \reg{AX}, dan isi \eng{queue}.
        \item Ulangi sampai instruksi \code{INT 21h}. Hentikan
              penelusuran sebelum layanan DOS dijalankan.
        \item Susun hasil pengamatan dalam tabel seperti
              Tabel~\ref{tab:trace-program}.
    \end{enumerate}

    \item \textbf{Menelusuri dengan DEBUG.}
    \begin{enumerate}
        \item Dari DOSBox, rakit program dengan
              \file{tasm siklus-dasar.asm}, lalu link dengan
              \file{tlink siklus-dasar.obj}.
        \item Jalankan \file{debug siklus-dasar.exe}.
        \item Ketik \file{r} untuk menampilkan register.
        \item Ketik \file{t} untuk menjalankan satu instruksi, lalu
              catat perubahan \reg{IP} dan \reg{AX}.
        \item Bandingkan hasil DEBUG dengan hasil emu8086.
    \end{enumerate}

    \item \textbf{Menentukan unit yang bekerja.}
    \begin{enumerate}
        \item Untuk setiap instruksi pada tabel, tandai apakah
              pengambilan instruksi dikerjakan BIU, eksekusi
              dikerjakan EU, atau keduanya.
        \item Jelaskan mengapa \code{MOV AX, 1234h} tidak memerlukan
              siklus bus data tambahan.
    \end{enumerate}

    \item \textbf{Modifikasi program.}
    \begin{enumerate}
        \item Tambahkan instruksi \code{MOV BX, AX} setelah
              \code{MOV AX, 1234h}.
        \item Rakit ulang, lalu telusuri kembali.
        \item Catat perubahan \reg{BX} dan panjang instruksi yang
              ditambahkan.
    \end{enumerate}

    \item \textbf{Menghitung alamat fisik.}
    \begin{enumerate}
        \item Untuk setiap baris tabel, hitung alamat fisik
              \reg{CS}:\reg{IP} dengan rumus
              $(\reg{CS} \times 16) + \reg{IP}$.
        \item Tuliskan hasilnya dalam kolom tambahan.
    \end{enumerate}
\end{enumerate}

\subsection{Hasil yang Dikumpulkan}

\begin{enumerate}
    \item Gambar diagram blok BIU/EU beserta catatan koreksi.
    \item Tabel pengamatan \reg{CS}, \reg{IP}, \reg{AX}, \reg{DS},
          dan isi \eng{queue} untuk setiap langkah.
    \item Tabel alamat fisik yang dihitung dari pasangan
          \reg{CS}:\reg{IP}.
    \item Berkas \file{siklus-dasar.asm} versi modifikasi.
    \item Analisis singkat tentang perbedaan pekerjaan BIU dan EU.
\end{enumerate}

\clearpage
\begin{asesmenbox}
\vspace{1.1em}
\textbf{Asesmen praktikum 2 (bobot 5\% dari komponen praktikum).}

\begin{tabular}{@{}p{6.0cm}cp{4.8cm}@{}}
    \toprule
    \textbf{Indikator} & \textbf{Bobot} & \textbf{Kriteria} \\
    \midrule
    Ketepatan diagram blok BIU/EU & 30\% &
    Seluruh komponen utama berada pada unit yang benar \\
    \addlinespace
    Ketepatan tabel penelusuran & 30\% &
    Nilai \reg{IP}, \reg{AX}, dan isi antrean tercatat benar \\
    \addlinespace
    Ketepatan perhitungan alamat fisik & 25\% &
    Seluruh alamat dihitung dengan rumus yang benar \\
    \addlinespace
    Analisis pembagian tugas & 15\% &
    Perbedaan kerja BIU dan EU dijelaskan dengan tepat \\
    \bottomrule
\end{tabular}
\end{asesmenbox}
